\section{The rational essential image of extension of the finite field}

The results here are concrete consequences of the classical
Milne classification applied in CC-AG08, within the established
abelian, Artin, and Tate numerical category.  They are not
claims of new literature results.  In contrast to the split
coefficient statement CC-AG07, they retain rational
coefficients and actual objects, including their positive
multiplicities.

Fix \(q=p^a\).  Let the target field \(B\) be either
\(\mathbb F_{q^m}\), for a positive integer \(m\), or
\(\overline{\mathbb F}_p\).  In the notation of CC-AG08,
extension of the base field acts on split labels by the
surjective, \(\Gamma\)-equivariant map
\[
 f_B:W(q)\longrightarrow\Lambda_B,\qquad
 \alpha\longmapsto
 \begin{cases}
       \alpha^m,& B=\mathbb F_{q^m},\\
       [\alpha,a],& B=\overline{\mathbb F}_p.
 \end{cases}                                                   \tag{CC-AG.34}
\]
Here \(\Gamma=\operatorname{Gal}(\overline{\mathbb Q}/\mathbb Q)\).
For a \(\Gamma\)-orbit \(O\) in the source or target, write
\(S_O\) for its rational simple object and \(d_O\) for its
Schur index, with the exact local-invariant convention of
CC-AG08.

\subsection{CC-AG09: the orbit and index criterion for the essential image}

For a target orbit \(O'\), define
\[
 \mathscr L_B(O')=
   \{O\subset W(q): O\text{ is a }\Gamma\text{-orbit},
                         \ f_B(O)=O'\}.
\]
For each such source orbit, put
\[
 c_O=\frac{|O|}{|O'|},\qquad
 r_O=\frac{c_Od_O}{d_{O'}},\qquad
 H_B(O')=\left\{\sum_{O\in\mathscr L_B(O')}a_Or_O:
       a_O\in\mathbb Z_{\ge0},\ \text{finite support}\right\}.
                                                               \tag{CC-AG.35}
\]
Then \(c_O\) and \(r_O\) are positive integers, and
\[
             (S_O)_B\simeq S_{O'}^{\oplus r_O}.                \tag{CC-AG.36}
\]
In particular, the formula includes the arithmetic divisibility
\(d_{O'}\mid c_Od_O\).

Let an actual rational target object have its unique simple
decomposition
\[
                       M=\bigoplus_{O'}S_{O'}^{\oplus b_{O'}}.
                                                               \tag{CC-AG.37}
\]
It belongs to the essential image of
\(\mathcal A_q\to\mathcal A_B\) if and only if
\[
                         b_{O'}\in H_B(O')
                       \quad\text{for every }O'.              \tag{CC-AG.38}
\]
The rational isomorphism classes of its preimages are exactly
the finite-support nonnegative integer solutions of
\[
                  \sum_{O\in\mathscr L_B(O')}a_Or_O=b_{O'}
                      \quad\text{for every }O'.               \tag{CC-AG.39}
\]
No unspecified object-descent condition remains in this
criterion.

For a finite target extension, \(\mathscr L_B(O')\) is finite:
its source labels are precisely the roots of
\[
                      \prod_{\beta\in O'}(t^m-\beta)
                               \ \in\mathbb Q[t].             \tag{CC-AG.40}
\]
Thus the criterion consists of finitely many field and index
calculations and nonnegative integer equations for a given
target object.  For the geometric target the lift set can be
infinite, but each additive monoid \(H_B(O')\subset\mathbb N\)
is still finitely generated.  This latter assertion does
not supply an effective bound on the orders of roots of
unity that one must inspect.

On Grothendieck groups, the corresponding coefficient condition
is instead divisibility by
\(\gcd\{r_O:O\in\mathscr L_B(O')\}\).  That is the group
completion of (CC-AG.35), rather than a replacement for
the stated actual-object condition.

\paragraph{Complete proof.}
The image of the transitive finite \(\Gamma\)-set \(O\)
under a \(\Gamma\)-equivariant map is the transitive set
\(O'\).  Its fibres all have the same cardinality \(c_O\):
an element of \(\Gamma\) carrying one target label to another
gives a bijection between the corresponding fibres.
Consequently \(c_O=|O|/|O'|\) is a positive integer.
If \(\alpha\in O\), the stabilizer inclusion also identifies
\[
       c_O=[K_\alpha:K_{f_B(\alpha)}].
                                                               \tag{CC-AG.41}
\]
Here \(K_\alpha=\mathbb Q(\alpha)\); for a germ the target
field is its stabilized centre from (CC-AG.24).
The field inclusion is literal for the chosen labels,
because a sufficiently large power representing the germ
lies in \(\mathbb Q(\alpha)\).

After extending coefficients to \(E\), CC-AG08 gives
\((S_O)_E\) multiplicity \(d_O\) on each label of \(O\).
Extension of the base field merges exactly the labels
having the same image under \(f_B\).  Thus \((S_O)_{B,E}\)
has multiplicity \(c_Od_O\) on every label of \(O'\) and
zero elsewhere.
The rational classification in the target says that a
rational object with this support is a sum of copies of
the single simple \(S_{O'}\), whose split multiplicity is
\(d_{O'}\) at each of these labels.  Its number of copies is
therefore \(r_O=c_Od_O/d_{O'}\), a positive integer.
This proves (CC-AG.36) and the asserted index divisibility.

Every rational source object is a unique finite sum
\(\bigoplus_O S_O^{\oplus a_O}\).  Applying (CC-AG.36)
gives exactly (CC-AG.39).  This proves necessity of
(CC-AG.38).  Conversely, for each of the finitely many
target orbits with \(b_{O'}>0\), membership in the monoid
supplies a finite nonnegative expression for \(b_{O'}\).
Taking the corresponding finite sum of source simples
constructs a preimage.  Simples mapping to different
target orbits cannot interfere with these equations, and
every \(r_O\) is positive.  Thus a target coefficient zero
forces all the associated source coefficients to be zero.
Uniqueness of simple decompositions proves the preimage
classification (CC-AG.39), including the zero object.

For a finite extension, every eligible label is an
\(m\)-th root of a label of \(O'\), and every such root
is a Weil \(q\)-number by the root argument in CC-AG07.
The product (CC-AG.40) is Galois invariant and hence has
rational coefficients.  Its roots form a finite set,
partitioned by the source orbits, proving finiteness.
Their centre fields and the local-invariant formula of
CC-AG08 give the stated arithmetic data.

For the geometric target, surjectivity of (CC-AG.34)
ensures that the set of \(r_O\)'s is nonempty.  Let
\(c\) be its least member; it is also the least positive
element of \(H_B(O')\).  For each residue class modulo
\(c\) which meets \(H_B(O')\), let \(h_i\) be the least
member in that class.  Every \(h\in H_B(O')\) is of
the form \(h_i+tc\) with an integer \(t\ge0\).
Hence \(H_B(O')\) is generated by \(c\) and the finitely
many \(h_i\).  Each \(h_i\), by definition of the monoid,
is itself a finite sum of the original \(r_O\)'s.
Taking the union of these finitely many finite supports
shows that finitely many of the actual lift indices
already generate \(H_B(O')\).  The argument proves
existence, without an effective cutoff on the lift set.

Finally, allowing arbitrary integer, rather than
nonnegative integer, source multiplicities gives the
subgroup of \(\mathbb Z\) generated by the \(r_O\)'s.
That subgroup is their greatest common divisor times
\(\mathbb Z\).  Independence of the target orbit
components proves the Grothendieck-group assertion.
\(\square\)

\subsection{CC-AG10: an exact even-multiplicity obstruction over
\(\mathbb F_3\)}

Set
\[
 K=\mathbb Q(s),\quad s^2=-5,\qquad
 \beta=2+s\in W(9),\qquad
 \xi=[\beta,2]\in\mathcal W_3.
                                                               \tag{CC-AG.42}
\]
Let \(T_9\in\mathcal A_9\) and
\(T_{\mathrm{geom}}\in\mathcal A_{\overline{\mathbb F}_3}\)
be the simple rational objects of the orbits of
\(\beta\) and \(\xi\), respectively.  Both have rank two,
endomorphism field \(K\), and Schur index one, and
\((T_9)_{\overline{\mathbb F}_3}\simeq T_{\mathrm{geom}}\).

For every nonnegative integer \(b\), either one of the
objects
\[
                T_9^{\oplus b},\qquad
                T_{\mathrm{geom}}^{\oplus b}
\]
is the base extension of an object in \(\mathcal A_3\)
if and only if \(b\) is even.
Thus the geometric multiplicity monoid in (CC-AG.35)
for this orbit is exactly \(2\mathbb N\).  In particular,
extension of the base field
\(\mathcal A_3\to\mathcal A_{\overline{\mathbb F}_3}\)
is not essentially surjective over rational coefficients.

\paragraph{Complete proof.}
We first establish the small number-field facts explicitly.
The ring of integers of \(K\) is
\(\mathcal O_K=\mathbb Z[s]\).  Indeed, if \(x+ys\)
with \(x,y\in\mathbb Q\) is integral, its trace and norm
are integers, so \(2x\in\mathbb Z\) and
\(x^2+5y^2\in\mathbb Z\).  Their discriminant
\(-20y^2\) is an integer.  Writing \(y\) in lowest
terms shows that its denominator is at most two, because
its denominator squared divides \(20\).
Thus \(x=A/2,\ y=B/2\) for integers \(A,B\), and
\[
                         A^2+5B^2\equiv0\pmod4
\]
forces both \(A\) and \(B\) to be even.  This proves one
inclusion; the reverse inclusion is immediate because
\(s\) is integral.

The two ideals
\[
       \mathfrak p=(3,s-1),\qquad
       \overline{\mathfrak p}=(3,s+1)
\]
are distinct prime ideals with residue field \(\mathbb F_3\).
Reduction of \(s^2+5\) modulo \(3\) gives
\((s-1)(s+1)\) with distinct roots, and the Chinese
remainder theorem gives
\[
                     (3)=\mathfrak p\overline{\mathfrak p}.
\]
Now \(\beta\) vanishes modulo \(\mathfrak p\) and is
equal to \(1\) modulo \(\overline{\mathfrak p}\).
Its field norm is \(4+5=9\).
Unique ideal factorization and the multiplicativity of
the norm therefore give
\[
                           (\beta)=\mathfrak p^2.             \tag{CC-AG.43}
\]
The ideal \(\mathfrak p\) is not principal.  A generator
would be an integer \(x+ys\in\mathcal O_K\) satisfying
\[
                           x^2+5y^2=3,
\]
which has no integral solution: nonzero \(y\) makes the
left side at least five, and \(y=0\) would make three
a square.  Consequently the ideal class of
\(\mathfrak p\) has order exactly two.  The ideal facts
used here are Milne (2020), Theorem~3.20, p.~53, and
Propositions~4.1--4.2, p.~69; no class-number table is
being used.

The two complex conjugates of \(\beta\) have absolute
value \(3\), and (CC-AG.43) shows that it is a \(3\)-unit.
Hence \(\beta\) is a Weil \(9\)-number of weight one.
Every root of unity in \(K\) is an integral unit; its
positive norm is one, and the equation
\(x^2+5y^2=1\) shows that it is \(1\) or \(-1\).
The ratio \(\beta/\overline\beta\) is neither \(1\)
nor \(-1\), since \(\beta\) has both nonzero real and
nonzero imaginary parts.  It is therefore not a root
of unity.  It follows that
\(\mathbb Q(\beta^N)=K\) for every \(N\ge1\):
otherwise the two conjugates would have equal
\(N\)-th powers.  In particular, \(K_\xi=K\), and the
germ orbit has two elements.

The two places above \(3\) have local degree one.
At them the slopes of \(\beta\), relative to \(q_0=9\),
are respectively \(2/2=1\) and \(0/2=0\).
There are no real places.  Formula (CC-AG.25) therefore
makes every local invariant zero, both for \(\beta\)
and for its stabilized germ.  It follows that
\[
       d_\beta=d_\xi=1,\qquad
       D_\beta=D_\xi=K,\qquad
       \operatorname{rank}T_9
         =\operatorname{rank}T_{\mathrm{geom}}=2.
                                                               \tag{CC-AG.44}
\]
The label-to-germ map is a bijection on these two-element
orbits.  Equation (CC-AG.36) now gives the claimed
geometric base change of \(T_9\).

Consider any source orbit which maps to the orbit of
\(\xi\).  Replacing a chosen label \(\alpha\in W(3)\)
by a Galois conjugate if necessary, arrange that
\([\alpha,1]=\xi\).  By the definition of a germ,
for some integer \(N\ge1\),
\[
                            \alpha^{2N}=\beta^N.             \tag{CC-AG.45}
\]
Put \(L=\mathbb Q(\alpha)\).
Since \(\mathbb Q(\beta^N)=K\), this equation proves
the required inclusion \(K\subset L\).
It also proves that \(\alpha\) is integral: it is a
root of a monic polynomial with coefficient
\(\beta^N\), an algebraic integer.
Passing to fractional ideals in \(\mathcal O_L\),
(CC-AG.43)--(CC-AG.45) give
\[
                      (\alpha)^{2N}
                           =(\mathfrak p\mathcal O_L)^{2N}.
\]
The group of fractional ideals is free abelian on the
prime ideals, so its power maps are injective and
\[
                            (\alpha)=\mathfrak p\mathcal O_L.
                                                               \tag{CC-AG.46}
\]
The relative ideal norm sends this equality to
\[
             \bigl(N_{L/K}(\alpha)\bigr)
                        =\mathfrak p^{[L:K]}.                 \tag{CC-AG.47}
\]
This is exactly Proposition~4.1(a),(c) of Milne (2020),
p.~69; the extension \(L/K\) need not be Galois.
Since the class of \(\mathfrak p\) has order two,
(CC-AG.47) forces \([L:K]\) to be even.

For clarity, the source Schur index is also exactly one.
The field \(L\) contains the imaginary quadratic field
\(K\), so it has no real embeddings.  At a place \(v\)
above \(3\), let \(e_v,f_v\) be its ramification index
and residue degree.  The normalized integral valuation
satisfies
\(\operatorname{ord}_v(3)=e_v\) and
\([L_v:\mathbb Q_3]=e_vf_v\).
Since the source field is \(\mathbb F_3\), its invariant
from (CC-AG.25) is
\[
        -\frac{\operatorname{ord}_v(\alpha)}{e_v}
                          e_vf_v
          =-\operatorname{ord}_v(\alpha)f_v
          \equiv0\pmod{\mathbb Z}.
                                                               \tag{CC-AG.48}
\]
All other invariants vanish as well.  Hence \(d_\alpha=1\).
The transition multiplicity is therefore
\[
        r_O=\frac{[L:\mathbb Q]}{[K:\mathbb Q]}
                              \frac{d_\alpha}{d_\xi}
             =[L:K],
                                                               \tag{CC-AG.49}
\]
which is even.  Every rational source object is a sum
of such simples, proving that a geometric preimage can
exist only in even multiplicity.

Finally, choose \(\alpha\) with \(\alpha^2=\beta\).
It is a Weil \(3\)-number by the root calculation
of CC-AG07.  The field \(L=\mathbb Q(\alpha)\)
contains \(K\).  The element \(\beta\) is not a
square in \(K\): a square root in \(K\) would be
integral and would have positive norm three, contrary
to the equation already considered.  Thus
\([L:K]=2\).  Equivalently, the minimal polynomial
of \(\alpha\) is the irreducible polynomial
\[
                              t^4-4t^2+9.
\]
The index computation (CC-AG.48) applies, so \(d_\alpha=1\),
and its simple rational object has rank four.
Equation (CC-AG.36) gives two copies of
\(T_{\mathrm{geom}}\).  Finite extension to \(\mathbb F_9\)
already gives two copies of \(T_9\), since the power map
has orbit fibre size two and both relevant indices
are one.  Taking arbitrary finite sums realizes every
even multiplicity.
An odd-multiplicity finite-field preimage would give an
odd-multiplicity geometric preimage after further base
extension, which has just been excluded.
This proves both asserted equivalences and
\(H_{\overline{\mathbb F}_3}(O_\xi)=2\mathbb N\).
\(\square\)

\subsection{Scope and primary references}

The criterion describes the exact rational essential image
within \(\mathcal A_q\).  It does not assert that every
smooth projective variety belongs to that category, that
numerical projectors lift to Chow projectors, or that the
rational category has a neutral fibre functor.  The latter
has already been disproved by AQ-RATIONAL-FIBRE.  The
arbitrary-variety Tate gate remains the one separated in
CC-XC01--04 and BR-EFF01.

\begin{itemize}
\item Milne, J.~S. (1994). Motives over finite fields. In
U.~Jannsen, S.~Kleiman, \& J.-P.~Serre (Eds.),
\emph{Motives} (Proceedings of Symposia in Pure Mathematics,
Vol.~55, Part~1, pp.~401--459). American Mathematical
Society. Corrected author version dated November~2, 2019:
\url{https://www.jmilne.org/math/xnotes/Seattle1.pdf}.
The classification, rank, invariant, and scalar-extension
inputs are the exact statements listed in CC-AG08.
\item Milne, J.~S. (2020, July~19).
\emph{Algebraic number theory} (Version~3.08)
[Author course notes].
\url{https://www.jmilne.org/math/CourseNotes/ANT.pdf}.
Exact ideal inputs: Theorem~3.20, printed p.~53;
Proposition~4.1(a),(c) and Proposition~4.2(a),
printed p.~69.  These are ideal factorization and norm
identities, rather than a precomputed assertion about
the example's ideal class.
\end{itemize}
